\documentclass[12pt]{article}
\usepackage[utf8]{inputenc}
\usepackage[T1]{fontenc}
\usepackage[margin=1in]{geometry}
\usepackage{setspace}
\usepackage{xcolor}
\usepackage[colorlinks=true, urlcolor=blue!60!black, linkcolor=blue!60!black]{hyperref}
\setlength{\parindent}{0em}
\setlength{\parskip}{0.6em}
\renewcommand{\baselinestretch}{1.0}
\begin{document}
\begin{center}
{\LARGE \textit{Hein v. Freedom From Religion Foundation, Inc.}, 551 U.S. 587 (2007)\par}
\vspace{0.5em}
{\large Case Brief --- Two-Page Edition\par}
\end{center}

\section{Memory Jogger}
Taxpayers could not challenge White House faith-based conferences: \textit{Flast v.\ Cohen} covers congressional spending mandates, not discretionary executive spending (\href{https://supreme.justia.com/cases/federal/us/551/587/}{551 U.S. 587}).

\section{Detailed Case Facts}
President Bush created faith-based offices by executive order, funded from general lump-sum appropriations rather than any specific statute. The Freedom From Religion Foundation and members sued as taxpayers, alleging that conferences praising faith-based groups violated the Establishment Clause.

\section{Procedural History}
The district court dismissed for lack of standing; a divided Seventh Circuit, per Judge Posner, reversed; the Supreme Court reversed.

\section{Judicial Votes}
Reversed, 5--4. Alito wrote a plurality for Roberts and Kennedy; Kennedy concurred; Scalia, with Thomas, concurred in the judgment; Souter, with Stevens, Ginsburg, and Breyer, dissented.

\section{Holding}
\textit{Flast}'s taxpayer-standing exception requires a challenge to spending expressly authorized or mandated by a specific congressional enactment under the Taxing and Spending Clause; discretionary executive spending from general appropriations does not qualify.

\section{Analysis of Opinions}
\textbf{Plurality:} \textit{Flast} is confined to its facts; extending it would make courts ``continuing monitors'' of the Executive. \textbf{Kennedy:} judges must not become ``speech editors'' for executive officials. \textbf{Scalia:} the line is meaningless; \textit{Flast}'s ``Psychic Injury'' theory should be overruled. \textbf{Souter:} the taxpayer's injury is identical whichever branch spends the money.

\section{Examples: Future Applications}
\textbf{Same-Side:} agency workshops for Christian charities paid from a general salaries account. \textbf{Opposite-Side:} a statute appropriating \$50 million for grants to ``houses of worship and sectarian ministries.'' \textbf{Fence-Sitter:} a religion-neutral grant statute whose committee report earmarks funds for prison ministries.

\section{Critique}
Originalists call the plurality's distinction unprincipled ``stealth overruling'' that leaves an ``arbitrary'' doctrine. Progressive scholars such as Ira Lupu and Robert Tuttle say \textit{Hein} opens a large loophole in Establishment Clause enforcement for executive programs.

\section{Key Quotations}
\begin{quote}
``It is a necessary concomitant of the doctrine of \textit{stare decisis} that a precedent is not always expanded to the limit of its logic. That was the approach that then-Justice Rehnquist took in his opinion for the Court in \textit{Valley Forge}, and it is the approach we take here. We do not extend \textit{Flast}, but we also do not overrule it. We leave \textit{Flast} as we found it.''
\end{quote}
\noindent --- Justice Alito (plurality opinion), 551 U.S. at 615.

\end{document}
